
%%
%% Part 2 – Theoretical Foundations of Diffusion Models
%%
%% This chapter revisits the core theory behind denoising diffusion
%% probabilistic models (DDPMs) and related variants.  It begins with
%% a derivation of the closed‑form marginal distribution in the forward
%% process, explains why predicting noise is advantageous, analyses
%% the variational objective, discusses DDIM sampling and introduces
%% classifier‑free guidance.  Additional sections elaborate on the
%% mathematical framework, network architectures, training objectives
%% and sampling algorithms.  All content has been carefully preserved
%% and reorganised from the original assignment.

\section{Theoretical Foundations of Diffusion Models}

Diffusion models are a family of generative techniques that learn to
reverse a gradually applied corruption process.  By viewing
generation as the solution to a stochastic or deterministic
differential equation, these models are grounded in probability
theory and admit closed‑form expressions for certain quantities.
This section develops the key theoretical results and motivations.

\subsection{Closed‑Form Expression for the Forward Process}

Consider a Markov chain $\{x_t\}_{t=0}^T$ that progressively
corrupts a clean sample $x_0$ by adding Gaussian noise at each step.
The single‑step forward transition is defined as
\begin{equation}
  q(x_t \mid x_{t-1}) = \mathcal{N}(x_t; \sqrt{\alpha_t}\,x_{t-1}, (1-\alpha_t)I),
\end{equation}
where $\alpha_t \in (0,1]$ controls the variance schedule and $I$ is
the identity matrix.  Expanding this recursion yields a closed‑form
expression for the marginal at time $t$ given $x_0$.

\begin{solutionbox}
  To prove that $q(x_t\mid x_0) = \mathcal{N}(\sqrt{\bar{\alpha}_t} x_0,
  (1 - \bar{\alpha}_t)I)$, where $\bar{\alpha}_t = \prod_{s=1}^t
  \alpha_s$, we begin with the reparameterisation
  \begin{equation}
    x_t = \sqrt{\alpha_t}\,x_{t-1} + \sqrt{1-\alpha_t}\,\epsilon_{t-1},
    \quad \epsilon_{t-1} \sim \mathcal{N}(0,I).
  \end{equation}
  Substituting the expression for $x_{t-1}$ recursively back to
  $x_0$ accumulates products of square roots:
  \begin{equation}
    x_t = \sqrt{\alpha_t\alpha_{t-1}}\,x_{t-2} + \sqrt{\alpha_t(1-\alpha_{t-1})}
    \epsilon_{t-2} + \sqrt{1-\alpha_t}\,\epsilon_{t-1}.
  \end{equation}
  Merging independent Gaussian terms results in a single Gaussian
  noise term with variance $1 - \alpha_t\alpha_{t-1}$.  Induction
  shows that
  \begin{equation}
    x_t = \sqrt{\bar{\alpha}_t}\,x_0 + \sqrt{1 - \bar{\alpha}_t}\,\epsilon,
    \quad \epsilon \sim \mathcal{N}(0,I),
  \end{equation}
  which implies the stated distribution.
\end{solutionbox}

The above proof reveals that the forward process admits a closed‑form
solution for the marginal distribution at any timestep $t$.
Therefore, the forward diffusion can be sampled directly without
iterating through all intermediate states.

\subsection{Why Predict Noise?}

The reverse diffusion process aims to reconstruct $x_{t-1}$ from
$x_t$.  One could train a neural network to predict the posterior
mean $\mu_\theta(x_t, t)$ directly.  However, it is more stable to
predict the noise component $\epsilon$ added during the forward step.
Two key advantages of noise prediction are summarised below.

\begin{solutionbox}
  \begin{enumerate}[label=\arabic*., leftmargin=2em]
    \item \textbf{Scale invariance and numerical stability.}  The
      target noise $\epsilon \sim \mathcal{N}(0,I)$ has constant
      variance across timesteps.  Predicting $\mu_\theta$ or $x_0$ would
      require dealing with widely varying scales: at early timesteps
      the data dominates, while at late timesteps the noise dominates.
      Noise prediction decouples the signal magnitude from the noise
      level, leading to more stable gradients.
    \item \textbf{Improved sample quality via reweighted loss.}
      Minimising the variational lower bound (VLB) yields a loss with
      timestep‑dependent weights that emphasise easy early steps and
      down‑weight difficult later steps.  By predicting noise and using
      the simple MSE loss $\|\epsilon - \epsilon_\theta(x_t,t)\|^2$,
      one effectively reweights the VLB, encouraging the model to
      focus on the most challenging denoising tasks.  Empirically
      this produces sharper and more realistic samples.
  \end{enumerate}
\end{solutionbox}

\subsection{Variational Lower Bound and KL Terms}

Training diffusion models can be viewed as maximising a variational
lower bound (VLB) on the data log‑likelihood.  The VLB decomposes
into three terms: a reconstruction term $L_0$, a sum of KL
divergences $L_{1:T-1}$ that penalise mismatch between the learned
reverse transitions and the true posteriors, and a prior matching
term $L_T$.  The KL terms encourage the reverse process to match the
forward process, ensuring that the model learns to invert the
corruption correctly.  In practice these terms are ignored in
favour of the simplified noise‑prediction loss; see the following box
for details.

\begin{solutionbox}
  \textbf{(A) Role of KL terms.}  The KL divergences
  \(D_{\text{KL}}\bigl(q(x_{t-1}\mid x_t, x_0)\,\|\,p_\theta(x_{t-1}\mid x_t)\bigr)\)
  encourage the neural network to produce reverse transitions that
  match the true forward posteriors.  This ensures that the learned
  reverse diffusion is consistent with the forward corruption process.

  \textbf{(B) Why ignore them?}  In practice the KL terms simplify to
  a weighted squared error between the predicted noise and the true
  noise.  The weights decrease for small $t$, rendering early terms
  nearly negligible.  Ignoring the KL weights yields the simplified
  MSE loss, which focuses learning on the more difficult denoising
  steps and often improves perceptual quality despite a slight loss
  in likelihood.
\end{solutionbox}

\subsection{DDIM Sampling}

\emph{Denoising Diffusion Implicit Models} (DDIM) introduce a
Deterministic sampling rule derived from the reverse SDE of diffusion.
In the DDPM sampling rule, a small amount of noise $\eta$ is added
at each step.  DDIM sets this noise scale to zero ($\eta=0$), making
The sampling process deterministic and non‑Markovian.  The update
Equation becomes
\begin{equation}
  x_{t-1} = \sqrt{\alpha_{t-1}}
  \Bigg(
    \frac{x_t - \sqrt{1-\bar{\alpha}_t}\,\epsilon_\theta(x_t,t)}{\sqrt{\bar{\alpha}_t}} +
    \sqrt{1-\frac{\bar{\alpha}_t}{\bar{\alpha}_{t-1}}}\,\epsilon_\theta(x_t,t)
  \Bigg).
\end{equation}
By skipping over timesteps (e.g. using 50 steps instead of 1000), one
can sample much more quickly at minimal quality loss.  The
following box summarises the mechanism and benefits.

\begin{solutionbox}
  \textbf{(A) How does DDIM make sampling deterministic?}  By
  setting the noise scale $\eta$ to zero in the reverse SDE, the
  random noise term disappears from the update rule.  As a result,
  the mapping from $x_t$ to $x_{t-1}$ becomes a deterministic
  function of $x_t$ and the predicted noise $\epsilon_\theta(x_t,t)$.
  This transforms the random walk of DDPM into a deterministic
  trajectory analogous to solving an ODE.

  \textbf{(B) Why is it faster?}  Because the reverse process is
  non‑Markovian, one can jump directly from timestep $t$ to $t-k$ by
  chaining together the deterministic update rule.  In practice, one
  selects a subset of timesteps (e.g. 50 out of 1000) and performs
  updates only at those points.  This reduces computational cost by a
  factor of twenty while preserving fidelity for moderate noise levels.
\end{solutionbox}

\subsection{Classifier‑Free Guidance}

Guided generation traditionally uses a separate classifier to steer
unconditional models toward desired attributes.  \emph{Classifier‑free
Guidance} (CFG) eliminates the need for an external classifier by
training a single model on both conditional and unconditional inputs.
During training, conditioning information (e.g. a text prompt) is
randomly dropped and replaced with a null token.  At inference time,
the final noise prediction is a linear combination of the
conditional prediction and the unconditional prediction:
\begin{equation}
  \tilde{\epsilon}_\theta(x_t,c) =
  \epsilon_\theta(x_t,\emptyset) + w \bigl(
    \epsilon_\theta(x_t,c) - \epsilon_\theta(x_t,\emptyset)
  \bigr).
\end{equation}
Here $w$ is the guidance scale controlling the trade‑off between
adherence to the conditioning and output diversity.  The next box
summarises the idea.

\begin{solutionbox}
  \textbf{(A) Main idea of CFG.}  The model is trained to produce
  both unconditional predictions $\epsilon_\theta(x_t,\emptyset)$ and
  conditional predictions $\epsilon_\theta(x_t,c)$.  At inference,
  these are combined: starting from the unconditional prediction, the
  difference between the conditional and unconditional outputs is
  scaled by $w$ and added back.  This mimics a classifier gradient
  while requiring no separate classifier.

  \textbf{(B) Effect of guidance scale.}  Increasing $w$ produces
  images that adhere more closely to the conditioning (e.g. text
  prompts) but at the expense of diversity.  Too small a value leads
  to generic or off‑prompt outputs; too large a value results in
  oversaturation, unnatural textures and diminished variation.  An
  intermediate choice (e.g. $w\in[5,10]$ for images) often yields
  the best balance between fidelity and diversity.
\end{solutionbox}

\subsection{Mathematical Foundations}

Having introduced the key theoretical building blocks, we summarise
The mathematical framework underpinning diffusion models.  The
forward diffusion is a discrete‑time Markov chain defined by a
variance schedule $\beta_1,\dots,\beta_T$.  Its marginal distribution
admits a closed form as derived above.  The reverse process is
parametrised by a neural network that learns to denoise.  Training
amounts to minimising the expected squared error between predicted
and true noise or, equivalently, maximising a variational lower
Bound.  The probability flow ODE associated with the diffusion SDE
provides a deterministic counterpart used in DDIM.

\subsection{U‑Net Architecture for Diffusion Models}

Although the theory of diffusion models is agnostic to the choice of
neural architecture, U‑Nets have emerged as a powerful backbone for
noise prediction.  The U‑Net used in this report features four
Downsampling levels and four upsampling levels.  Each level contains
residual blocks with group normalisation and SiLU activation,
self‑attention modules, and 1×1 or transpose convolutional layers to
adjust channel dimensions.  Time embeddings generated via
sinusoidal encodings are injected into every residual block via
learned affine transformations.  Skip connections concatenate
feature maps from the encoder to the corresponding decoder blocks.

\subsection{Training Objective}

The simplified objective used throughout this report is
\begin{equation}
  \mathcal{L}_{\text{simple}}(\theta) =
  \mathbb{E}_{x_0,\epsilon,t}\bigl[\|\epsilon -
    \epsilon_\theta(x_t,t)\|_2^2\bigr].
\end{equation}
Here $t$ is sampled uniformly from $\{0,\dots,T-1\}$ for each
training example, $\epsilon \sim \mathcal{N}(0,I)$ and $x_t$ is
Generated from $x_0$ via the forward diffusion.  This loss focuses
training on accurately predicting noise across all timesteps and
implicitly emphasises harder denoising tasks.

\subsection{Sampling Algorithms}

Sampling from a trained diffusion model begins with pure Gaussian
noise $x_T \sim \mathcal{N}(0,I)$ and applies the learned reverse
transition repeatedly until reaching $x_0$.  The DDPM update rule
uses the posterior mean and variance of the reverse transition
Together with a noise term.  DDIM modifies this rule to remove
stochasticity and enable step skipping.  The pseudo‑code for DDPM
sampling is given below:
\begin{lstlisting}[language=Python, caption=DDPM Sampling Pseudo‑code]
for t in range(T, 0, -1):
    z = torch.randn_like(x_t) if t > 1 else 0
    eps = epsilon_theta(x_t, t)
    x_{t-1} = (1/sqrt(alpha_t)) * (x_t - (1-alpha_t)/sqrt(1-alpha_bar_t) * eps) + sigma_t * z
\end{lstlisting}

For DDIM sampling, one constructs a custom timetable of timesteps
$\{t_0=T, t_1, \dots, t_K=0\}$ (e.g. using a strided schedule) and
applies the deterministic update derived earlier.

\subsection{Conclusion of Part~2}

This chapter has provided a thorough theoretical overview of
Denoising diffusion models.  In addition to deriving the closed
Forward distribution, we have argued for noise prediction, analysed
The VLB, explained DDIM sampling and classifier‑free guidance, and
Discussed architectural and training considerations.  The next
chapter (Part~3) will implement these ideas in practice.
